\documentclass[12pt]{article}
\usepackage{amsmath}
\usepackage{amsfonts}
\usepackage{amssymb}
\usepackage{physics}
\usepackage{geometry}
\usepackage{parskip}
\geometry{a4paper, margin=0.7in}


\title{02YMECH - Solution 9}
\date{Assigned for the week of Nov 17, 2025}
\author{}

\begin{document}
\maketitle

\section*{Solutions}

\begin{enumerate}

\item Integrating the velocities,
\[
\mathbf r_1(t)=(1+t,\tfrac12 t^2),\quad
\mathbf r_2(t)=(-1-t,\tfrac12 t^2),\quad
\mathbf r_3(t)=(0,\tfrac12 t^2).
\]

\begin{enumerate}
\item[(a)] Center of mass:
\[
\mathbf R_{\rm CM}(t)=\frac{1}{3}\sum_{i=1}^3 \mathbf r_i(t)
=(0,\tfrac12 t^2).
\]
Velocity and acceleration:
\[
\mathbf V_{\rm CM}(t)=(0,t),\qquad
\mathbf A_{\rm CM}=(0,1).
\]
The center of mass does \emph{not} move with constant velocity.

\item[(b)] Total linear momentum:
\[
\mathbf P=\sum m\mathbf v_i
=(1-1+0,\; t+t+t)=(0,3t).
\]
Since \(\dot{\mathbf P}=(0,3)\neq 0\), linear momentum is \emph{not conserved}.

\item[(c)] Total angular momentum about the origin:
\[
L_z=\sum m(x_i v_{iy}-y_i v_{ix}).
\]
For each particle,
\[
L_{1z}=t,\quad L_{2z}=t,\quad L_{3z}=0,
\]
so
\[
L_z=2t.
\]
Since \(dL_z/dt=2\neq0\), angular momentum is \emph{not conserved}.
\end{enumerate}

\item
\begin{enumerate}
\item[(a)] Total mass:
\[
M=\int_0^a\!\!\int_0^b \sigma_0\frac{x}{a}\,dy\,dx
=\frac{\sigma_0 ab}{2}.
\]

\item[(b)] Center of mass:
\[
x_{\rm cm}=\frac{1}{M}\int_0^a\!\!\int_0^b x\,\sigma_0\frac{x}{a}\,dy\,dx
=\frac{2a}{3},
\]
\[
y_{\rm cm}=\frac{1}{M}\int_0^a\!\!\int_0^b y\,\sigma_0\frac{x}{a}\,dy\,dx
=\frac{b}{2}.
\]

Thus
\[
(x_{\rm cm},y_{\rm cm})=\left(\frac{2a}{3},\frac{b}{2}\right).
\]
\end{enumerate}

\end{enumerate}

\end{document}
